\documentclass[10pt,a4paper]{amsart}
\usepackage[margin=19mm]{geometry}
\usepackage{amsmath,amssymb,mathtools}
\usepackage[scr=rsfs,cal=euler]{mathalfa}
\usepackage{xfrac}
\usepackage[unicode,hidelinks,bookmarks=false]{hyperref}
\newcommand{\ie}{i.e.\ }
\newcommand{\cf}{cf.\ }
\newcommand{\resp}{respectively\ }
\newcommand{\Subsection}{\subsection}
\providecommand{\nobreakdash}{\nobreak\hbox{-}}
\setlength{\emergencystretch}{2em}
\multlinegap0pt
% ---------------------------------------------------------------------------
% Editorial AMS wrapper prelude. The theorem environments and the mathematical macros below are
% transcribed from the paper's own preamble (cached-originals/source0.tex lines 1--83); the AMS amsart
% class, the named format packages of the pinned TeX Live runtime and this editorial formatting are not
% part of the author's text. No mathematics is changed.
% ---------------------------------------------------------------------------
\usepackage{algorithm}
\numberwithin{equation}{section}
\newtheorem{lemma}{Lemma}[section]
\newtheorem{question}[lemma]{Question}
\newtheorem{theorem}[lemma]{Theorem}
\newtheorem{proposition}[lemma]{Proposition}
\newtheorem{definition}[lemma]{Definition}
\newtheorem{corollary}[lemma]{Corollary}
\newtheorem{reduction}[lemma]{Reduction}
\newtheorem{hypothesis}[lemma]{Hypothesis}
\newtheorem{notation}[lemma]{Notation}
\newtheorem{conjecture}[lemma]{Conjecture}
\theoremstyle{definition}
\newtheorem{remark}[lemma]{Remark}
\newtheorem{example}[lemma]{Example}
\newcommand{\Gal}{\mathrm{Gal}}
\newcommand{\Hol}{\mathrm{Hol}}
\newcommand{\Inn}{\mathrm{Inn}}
\newcommand{\InHol}{\mathrm{InHol}}
\newcommand{\Aut}{\mathrm{Aut}}
\newcommand{\Out}{\mathrm{Out}}
\newcommand{\End}{\mathrm{End}}
\newcommand{\Sz}{\mathrm{Sz}}
\newcommand{\PP}{\mathcal{P}}
\newcommand{\Aff}{\mathrm{Aff}}
\newcommand{\Proj}{\mathbb{P}}
\newcommand{\GL}{\mathrm{GL}}
\newcommand{\F}{\mathbb{F}}
\newcommand{\Z}{\mathbb{Z}}
\newcommand{\NN}{\mathcal{N}}
\newcommand{\M}{\mathcal{M}}
\newcommand{\ord}{\mathrm{ord}}
\newcommand{\id}{\mathrm{id}}
\newcommand{\Sym}{\mathrm{Sym}}
\newcommand{\GG}{\mathcal{G}}
\newcommand{\TT}{\mathcal{T}}
\newcommand{\bg}{\mathbf{g}}
\newcommand{\bh}{\mathbf{h}}
\newcommand{\bt}{\mathbf{t}}
\newcommand{\bu}{\mathbf{u}}
\newcommand{\bv}{\mathbf{v}}
\newcommand{\bw}{\mathbf{w}}
\newcommand{\bx}{\mathbf{x}}
\newcommand{\hei}{\mathrm{ht}}
\newcommand{\N}{\mathbb{N}}
\newcommand{\LRA}{\Leftrightarrow}
\newcommand{\conj}{\mathrm{conj}}
\newcommand{\opp}{\mathrm{op}}
\newcommand{\inv}{\mathrm{inv}}
\newcommand{\tkap}{\widetilde{\kappa}}
\newcommand{\Soc}{\mathrm{Soc}}
\newcommand{\Mat}{\mathrm{Mat}}
\newcommand{\Autsb}{\mathrm{Aut}_{\mathrm{sb}}}
\newcommand{\Autgp}{\mathrm{Aut}_{\mathrm{gp}}}
\newcommand{\Autmag}{\mathrm{Aut}_\mathrm{mag}}
\newcommand{\Symo}{\mathrm{Sym}_0}
\newcommand{\ttau}{\tilde{\tau}}
\newcommand{\tsigma}{\tilde{\sigma}}
\newcommand{\tzeta}{\tilde{\zeta}}
\newcommand{\tpi}{\tilde{\pi}}
\newcommand{\bs}{\backslash}
\newcommand{\Fix}{\mathrm{Fix}}
\newcommand{\Supp}{\mathrm{Supp}}

\begin{document}
\title{Magma Automorphisms and Quasi-Linear Cycle Sets}
\author{Nigel P. Byott, Edgar Jasko}
\date{Source: arXiv:2609.20028v1; distributed under CC BY 4.0}
\maketitle
\noindent\emph{Source, version and licence.} \emph{Magma Automorphisms and Quasi-Linear Cycle Sets}, by Nigel P. Byott, Edgar Jasko. The source of this editable unit is the e-print of \textbf{version v1} of arXiv:2609.20028, https://arxiv.org/abs/2609.20028v1, whose material is distributed under the Creative Commons Attribution 4.0 International licence, http://creativecommons.org/licenses/by/4.0/, as recorded in the arXiv licence metadata for this paper and shown on that abstract page. The whole of the body below is version v1, the newest version of the paper. Full rights notice: licenses/arxiv-2609.20028-CC-BY-4.0.txt.\par\smallskip
\noindent\textit{Editorial scope.} Counted unit: original Theorem 3.2 of the article, the statement carried by the source environment \texttt{theorem} with source label \texttt{trans-thm}, cached-originals/source0.tex lines 550--558, printed as \textbf{Theorem 3.2} exactly as in the article: for the cyclic group of order n greater than two and a transposition tau = (ab) with gcd(b-a,n) = 1, the group of magma automorphisms fixing tau is trivial when b is not equal to minus a and equals the two-element group generated by the inversion when b equals minus a, which also settles the conjecture of the paper in that case. It is delivered together with the authors' own complete proof as the single primary proof body, source0.tex lines 559--618 (60 lines): the \texttt{proof} environment that belongs to the statement and establishes it in full. No proof marker is added and no mathematics is written, repaired or re-derived; the authors' notation and the printed slips of the source are kept literal. Necessary complete context is retained uncounted, in source order: lines 197--199, the source \texttt{conjecture} with source label \texttt{main-conj} of the article that the retained passages refer to; lines 362--365, the source \texttt{proposition} with source label \texttt{autom} of the article that the retained passages refer to; lines 523--528, the source \texttt{proposition} with source label \texttt{S-lem} of the article that the retained passages refer to; lines 550--618, the source \texttt{theorem} with source label \texttt{trans-thm} of the article that the retained passages refer to. The article's own numbering is reproduced by explicit counter settings: the theorem-like counter of the article is set before each retained passage, so the counted statement prints with the number the article gives it. No \texttt{\textbackslash cite} command occurs in any retained passage: the closure of the retained passages over references and citations was computed, it contains no citation at all, so this unit delivers no bibliography entry and the article's own bibliography data is neither spliced into the bundled source nor redistributed. The source's own class is \texttt{amsart}; the e-print ships no class or style file, so no publisher class is shipped, used or redistributed, and the wrapper is the AMS \texttt{amsart} class with the article's own macros and theorem declarations transcribed from its preamble. The retained passages contain no figure environment and no graphics-inclusion command, so no artwork is redistributed. Every declared body of this unit is an exact, unmodified span of source0.tex: no body is a bounded passage comparison, \texttt{omit} was not used anywhere, and consequently no fidelity receipt is asserted in delivery-evidence.json. The complete concatenated original source is bundled as sources/210828/source0.tex. Omitted from this editable unit and therefore absent from it are source0.tex lines 1--196; 200--361; 366--522; 529--549; 619--1216. The AMS \texttt{amsart} wrapper, the named format packages of the pinned TeX Live runtime and this editorial source, licence and scope paragraph are editorial formatting only.\par\smallskip
\setcounter{section}{1}
\setcounter{equation}{8}
\setcounter{lemma}{1}
\begin{conjecture}  \label{main-conj}
Let $A$ be a finite abelian group, and let $\tau$, $\sigma \in \Symo(A)$. If $\sigma$ is a magma automorphism of $(A,\tau)$ then either $\Soc(A,\tau)\neq \{0\}$ or $\Soc(A,\sigma)\neq \{0\}$.
\end{conjecture}

\setcounter{section}{2}
\setcounter{equation}{2}
\setcounter{lemma}{5}
\begin{proposition} \label{autom}
If $\tau \in \Autgp(A)$ then 
$$ \Autmag(A,\tau) = \{ \sigma \in \Symo(A) : \sigma \tau = \tau \sigma \}. $$
\end{proposition}

\setcounter{section}{3}
\setcounter{equation}{0}
\setcounter{lemma}{0}
\begin{proposition} \label{S-lem}
Let $A$ be an arbitrary abelian group and let $\tau \in \Symo(A)$. Let
$$ S= \Supp(\ttau) = \{ a \in A : -\tau(-a) \neq a \},  $$
and let $\sigma \in \Autmag(A,\tau)$. Then, for each $b \in A$, we have  
	$$ \sigma(b+S) \backslash S =(\sigma(b)+S) \backslash S. $$
\end{proposition}

\setcounter{section}{3}
\setcounter{equation}{0}
\setcounter{lemma}{1}
\begin{theorem} \label{trans-thm}
Let $A=C_n$ with $n>2$, and let $\tau \in \Symo(A)$ be a transposition $\tau=(ab)$ where $\gcd(b-a,n)=1$. Then	
\begin{equation} \label{transp-aut}
		\Autmag(A,\tau) = \begin{cases} \{\id\} & \mbox{if } b \neq -a; \\
		\{\id, \inv\} & \mbox{if } b =-a;
	\end{cases} 
\end{equation}
where $\inv(x)=-x$. In particular, Conjecture \ref{main-conj} holds for $(A, \tau)$. 
\end{theorem}

\setcounter{section}{3}
\setcounter{equation}{1}
\setcounter{lemma}{2}
\begin{proof}

We begin with a preliminary observation. If $b=2a$ and $a=2b$ in $A$ then $3(b-a)=0$ and the coprimality hypothesis forces $n=3$. For $n=3$, we necessarily have $\tau=(12)=\inv$ and $\Autmag(A,\tau)=\{\id, \tau \}=\Autgp(A,\tau)$. Thus the Theorem holds for $n=3$. We henceforth assume that $n>3$, and further (by swapping $a$ and $b$ if necessary) that $b \neq 2a$. 

In the notation of Proposition \ref{S-lem}, we have 
$$ S= \Supp(\ttau)=	\{-a,-b\}. $$ 
For $0 \leq j \leq n-1$, let $T_j=a+j(a-b)+S = \{j(a-b), (j+1)(a-b)\}$. In particular, 
$$  	T_0 =  \{ 0, a-b\},  \qquad 	T_{n-1} = \{ b-a, 0 \}.  $$ 
The sets $T_j$ are pairwise distinct by the coprimality hypothesis, and for each $c \neq 0$ in $A$ there is a unique index $i$ such that 
$$ c \in T_j \LRA j=i \mbox{ or } i+1. $$ 
Moreover, there is a unique $j_0$ with $1 \leq j_0 \leq n-2$ such that $T_{j_0} = S= \{-a,-b\}$. We then have $T_{j_0-1} = \{b-2a, -a\}$ and $T_{j_0+1} = \{ -b, a-2b \}$. 

Next consider the sets $T'_j=T_j \bs S$. We have $T'_{j_0}= \emptyset$ and 
$$ T'_{j_0-1}=\{b-2a\}=\{(j_0-1)(a-b)\}, \quad T_{j_0+1} = \{(j_0+2)(a-b)\} = \{a-2b \}, $$ 
while $T'_j=T_j$ if $j \neq j_0$, $j_0 \pm 1$. 

Now let $\sigma \in \Autmag(A,\tau)$. By Proposition \ref{S-lem}, $\sigma$ permutes the sets $T'_j$. In particular, as $\sigma(0)=0$, the two sets $T'_0$, $T'_{n-1}$ containing $0$ must either be fixed by $\sigma$ or swapped by $\sigma$. 
We consider these two possibilities separately.
	
First suppose that $\sigma(T'_0)=T'_0$ and $\sigma(T'_{n-1}) = T'_{n-1}$. We will show that in this case $\sigma=\id$.
We first show by induction that for $0 \leq k \leq j_0-1$ we have 
\begin{equation} \label{transp-ind}
	\sigma(T'_k)=T'_k, \qquad \sigma(k(a-b))=k(a-b).
\end{equation}
This is true for $k=0$. Suppose that (\ref{transp-ind}) holds for some $k<j_0-1$. Then $T'_k=\{k(a-b),(k+1)(a-b)\}$ is fixed by $\sigma$. Since $\sigma$ fixes $k(a-b)$, it must also fix $(k+1)(a-b)$. Since $(k+1)(a-b) \in T'_j$ if and only if $j=k$ or $k+1$, and $\sigma$ fixes $T'_k$, it must also fix $T'_{k+1}$. Hence (\ref{transp-ind}) also holds for $k+1$, completing the induction. A similar argument, using descending induction and starting from $\sigma(T'_{n-1})=T'_{n-1}$, shows that (\ref{transp-ind}) also holds for $n-1 \geq k \geq j_0+2$. Thus $\sigma(c)=c$ for all $c \in A$ with the possible exceptions of $c=j_0(b-a)=-a$ and $c=(j_0+1)(b-a)=-b$. Hence $\sigma=\id$ or $\sigma=(-a,-b)=\ttau$. We will show that $\ttau$ cannot be a magma automorphism when $n>3$. Let $x=-2a$ and $y=b-2a$. Since $a \neq b$ and we have assumed that $b \neq 2a$, we have $\tau(-x)=-x$ and $\ttau(x)=x$. As $n>2$ and $\gcd(b-a,n)=1$, we have $2(b-a) \neq 0$, so that $b-2a \neq -a$, $-b$. Thus $\ttau(y)=y$. We now calculate
\begin{eqnarray*}
	 \ttau( x \cdot_\tau y) & = & \ttau(\tau(y-x)-\tau(-x)) \\
	                         & = & \ttau(\tau(b)-\tau(-x)) \\
	                         & = & \ttau(a-2a) \\
	                         & = & -b
\end{eqnarray*}
and
\begin{eqnarray*}
	\ttau(x) \cdot_\tau \ttau(y) & = & \tau( \ttau(y)-\ttau(x)) - \tau( -\ttau(x)) \\ 
								& = & \tau(y-x) - \tau(-x) \\
								& = & \tau(b)-\tau(-x) \\
								& = & a -(-x) \\
								& = & -a.						
\end{eqnarray*}        
As $-b \neq -a$ in $A$, this shows that $\ttau \not \in \Autmag(A,\tau)$.

Now suppose that $T'_0$ and $T'_{n-1}$ are swapped by the magma automorphism $\sigma$. We will show that in this case the only possibility is $\sigma =\inv$, and this can only occur if $b=-a$. 
We first show by induction that for $0 \leq k \leq j_0-1$, the following hold:
\begin{equation} \label{inv-bound}
	n-1-k \geq j_0+1; 
\end{equation}
\begin{equation} \label{inv-ind}
	\sigma(k(a-b))=k(b-a), \qquad  	\sigma(k(b-a))=k(a-b).
\end{equation}
\begin{equation} \label{inv-set}
	\sigma(T'_k) = T'_{n-1-k}, \qquad \sigma(T'_{n-1-k})=T'_k; 
\end{equation}
Clearly, these all hold when $k=0$. Suppose that they hold for some $k < j_0-1$. Then $|T'_k|=2$ and $T'_k=\{k(a-b), (k+1)(a-b)\}$. By (\ref{inv-set}), $\sigma$ swaps $T'_k$ and $T'_{n-1-k}$. Hence $|T'_{n-1-k}|=2$ as well, and $T'_{n-1-k}=\{k(b-a),(k+1)(b-a)\}$. In particular, $n-1-k \neq j_0+1$. Together with (\ref{inv-bound}), this shows that $n-1-k \geq j_0+2$, so that (\ref{inv-bound}) holds for $k+1$. Moreover, it follows from (\ref{inv-ind}) that $\sigma$ swaps $(k+1)(a-b)$ and $(k+1)(b-a)$, so (\ref{inv-ind}) holds for $k+1$. Finally, since 
$$  (k+1)(a-b) \in T'_j \LRA j=k \mbox{ or } j=k+1 $$
and
$$  (k+1)(b-a) \in T'_j \LRA j=n-1-k \mbox{ or } j= n-2-k, $$
it follows that $\sigma$ swaps $T'_{k+1}$ and $T'_{n-2-k}$. This shows that (\ref{inv-set}) holds for $k+1$, completing the induction. Now taking $k=j_0-1$ in (\ref{inv-bound}) and (\ref{inv-set}), we have that $n-j_0 \geq j_0+1$ and that $\sigma$ swaps $T'_{j_0-1}$ and $T'_{n-j_0}$. Thus $|T'_{n-j_0}|=|T'_{j_0-1}|=1$. Hence $n-j_0=j_0+1$, so that $n$ is odd and $j_0=\frac{1}{2}(n-1)$. As $T_{j_0}=\{j_0(a-b), (j_0+1)(a-b)\}=\{-a,-b\}$, we have $b=-a$. Since (\ref{inv-ind}) holds for $0 \leq k \leq \frac{1}{2}(n-1)$, we now know that $\sigma(c)=-c$ for all $c \in A$ with the possible exceptions of $c=\frac{1}{2}(n-1)(a-b)=-a=b$ and $c=\frac{1}{2}(n+1)(a-b)=-b=a$. Thus either $\sigma=\inv$ or $\sigma=\ttau \circ \inv$, where in this case $\ttau=\tau$. Now $\inv \in \Autgp(A)$ and $\inv$ commutes with $\tau$, so $\inv \in \Autmag(A,\tau)$ by Proposition \ref{autom}. We have already seen in the first case that $\ttau \not \in \Autmag(A,\tau)$. Thus $\ttau \circ \inv \not \in \Autmag(A,\tau)$. 

We have now shown both cases of (\ref{transp-aut}). In particular, $\Autmag(A,\tau) \subseteq \Autgp(A)$, so Conjecture \ref{main-conj} holds for $(A, \tau)$. 
\end{proof}

\end{document}
